\section[The Cooperation and Consistency Mechanism Pre-- and Post--Procedural Regulation]{The Cooperation and Consistency Mechanism\\Pre-- and Post--Procedural Regulation}

This chapter will outline the \gls{gdpr}'s cooperation and consistency mechanism.
It will first present the tasks and powers of \glspl{dpa} under the \gls{gdpr}. 
It will then recap the allocation of competences between the \gls{lsa} and the \glspl{sac}, followed by a summary of the pre-reform cooperation and consistency mechanism.
Finally, the chapter will analyze the key changes to the flow of the cooperation and consistency mechanism introduced by the Procedural Regulation in enforcement cases dealing with cross-border processing.

\subsection{The Tasks and Powers of Supervisory Authorities}

Under the \gls{gdpr}, \glspl{ms} are required to establish independent\autocite[52]{GDPR} \glspl{sa}, entrusted with monitoring the application of the Regulation.\autocite[51]{GDPR} 
Importantly, the \gls{gdpr} requires \glspl{sa} to all contribute to the application of the Regulation throughout the Union.\autocite[51(2)]{GDPR} 
The \gls{gdpr} fundamentally connects the institution of \glspl{sa} with the protection of the fundamental right to data protection, which is also the legal basis on which the Regulation was adopted.\autocites[preamble]{GDPR}[16(2)]{TFEU}

The \gls{sa} is competent for enforcing the Regulation within the territory of its \gls{ms},\autocite[51(1)]{GDPR} with well-defined tasks\autocite[57]{GDPR} and powers\autocite[58]{GDPR}.
Among the tasks of the \gls{sa} figure: the monitoring of the application of the \gls{gdpr};\autocite[57(1)(a)]{GDPR} the handling of complaints;\autocite[57(1)(f)]{GDPR} the cooperation with other \glspl{sa} to ensure the consistency of application of the \gls{gdpr};\autocite[57(1)(g)]{GDPR} and the handling of investigations, including with the input of other \glspl{sa}.\autocite[57(1)(h)]{GDPR}
The powers of the \gls{sa} include investigative, corrective, advisory powers, as well as the power to engage in legal proceedings.\autocite[58]{GDPR}

\subsection{Competence Allocation in Cross-Border Enforcement Cases}

In the case of cross-border processing, the \gls{lsa} becomes competent to supervise the processing activities of the controller or processor when its main or single establishment is in its territory.\autocite[56(1)]{GDPR}
The \gls{lsa} becomes the sole interlocutor of the controller or processor responsible for the cross-border processing at examination.\autocite[56(6)]{GDPR}
Cross-border processing refers to either the processing of personal data by a controller or processor with multiple establishments across the Union, or the processing of personal data by one entity that `affects or is likely to affect' data subjects in more than one \gls{ms}.\autocite[4(23)]{GDPR}

Other \glspl{sa} then may become `concerned'.\autocite[4(22)]{GDPR} 
An \gls{sac} is concerned when either: the processing in question is carried out by an establishment in its territory; the data subjects residing in the territory of its competence are affected by the processing; or when a complaint has been lodged with it.\autocite[4(22)(a)--(c)]{GDPR}

An exception to the general rule of competence of the \gls{lsa} is provided in Article 56.
When the subject of an investigation relates only to a specific \gls{ms}, the \gls{sac} of that \gls{ms} can be competent to carry out enforcement.\autocite[56(2)]{GDPR}
The \gls{sac} is however required to inform the \gls{lsa} of its willingness to start proceedings, and the \gls{lsa} can, within three weeks, take charge of the proceedings and conduct them based on the cooperation procedure.\autocite[56(3)]{GDPR}
However, in that case, the \gls{sac} has the possibility of submitting its own draft decision to the \gls{lsa}.\autocite[56(4)]{GDPR} 
The \gls{lsa} must have particular regard to the views of the \gls{sac} when it drafts its own decision.\autocite[56(4)]{GDPR}
If the \gls{lsa} does not wish to take over the proceedings, the \gls{sac} can proceed with the investigation and make use of the tools provided by the cooperation procedure.\autocite[56(5)]{GDPR}

\subsection{The Pre-Reform Cooperation and Consistency Mechanism}

\subsubsection{The Cooperation Procedure}

Article 60 and 61 \gls{gdpr} set out the cooperation procedure between the \gls{lsa} and the \glspl{sac} in cases involving cross-border processing. 
Fundamentally, the \gls{lsa} is required to cooperate with the \glspl{sac} to reach consensus, with the authorities exchanging `all relevant information with each other'.\autocite[60(1)]{GDPR}

While investigating potential infringements, \glspl{sa} can make reciprocal requests for mutual assistance.\autocite[60(2), 61(1)]{GDPR} 
Mutual assistance includes requests for information, investigations, and inspections.\autocite[61(1)]{GDPR}
Such a request must be answered within one month of its receipt, and the requested \gls{sa} must inform the requesting \gls{sa} of the outcome or progress on the request.\autocite[61(5)]{GDPR}
\glspl{lsa} can also make use of joint operations, where the \gls{lsa} and \glspl{sac} can carry out joint investigations, inspections, or monitor the application of corrective measures.\autocite[60(2), 62]{GDPR}
\glspl{sac} have a right to participate in joint operations with the \gls{lsa}.\autocite[62(2)]{GDPR}
In the case of joint operations, the host \gls{sa} may extend its investigative powers to the guest \gls{sa}, in accordance with national law.\autocite[62(3)]{GDPR}

When the \gls{lsa} has terminated its investigation, it must submit a draft decision to the \glspl{sac} and consult their views on the draft.\autocite[60(3)]{GDPR}
\glspl{sac} have four weeks to submit a `relevant and reasoned objection' to the draft.\autocite[60(4)]{GDPR}
An objection is regarded as `relevant and reasoned' if `clearly demonstrates the significance of the risks posed by the draft decision' concerning the fundamental rights and freedoms of data subjects.\autocite[4(24)]{GDPR}

An \gls{lsa} can follow the objections presented, and adjust the draft decision accordingly.\autocite[60(5)]{GDPR}
The revised decision will have to be submitted again to the \glspl{sac}, who will have two weeks to review it.\autocite[60(5)]{GDPR}
If the \gls{lsa} does not follow the objections, or it does not believe that they are relevant or reasoned, it must refer the matter to the consistency mechanism.\autocite[60(4)]{GDPR}

When the period of four weeks lapses, and no objections are raised, the \gls{lsa} can adopt the decision and the \glspl{sac} are deemed to have agreed to it.\autocite[60(6)]{GDPR} 
The \gls{lsa} then is required to inform the controller/processor of the adoption of the decision, as well as the \glspl{sac} and the \gls{edpb}.\autocite[60(7)]{GDPR}

When a decision is adopted, or a complaint is dismissed or rejected, the \glspl{sac} who received the original complaint are tasked with notifying the complainants of the decision.\autocite[60(7)--(8)]{GDPR}
When a complaint is only followed in part, the \gls{lsa} must inform the complainant of the points raised where it chose to act, while the \gls{sac} must adopt a separate decision dismissing the rejected parts of the complaint.\autocite[60(9)]{GDPR}

\subsubsection{The Consistency Mechanism}

The consistency mechanism is an extension of the cooperation procedure, meant to bring a consistent enforcement of the \gls{gdpr} across the Union.\autocites[63]{GDPR}[68]{googlevCNIL}[52]{facebookBelgium}
It is laid out in Articles 63--67 \gls{gdpr}, and can be triggered in various ways in the context of cross-border enforcement.

Firstly, the \gls{edpb} can issue non-binding opinions on matters of general application on the request of an \gls{sa}, its own chair, or the Commission.\autocite[64(2)]{GDPR}
This is especially relevant in case an \gls{sa} refuses to comply with a request for mutual assistance or joint operations.\autocite[64(2)]{GDPR}

Secondly, the \gls{edpb} can issue binding decisions in a number of instances.\autocite[65(1)]{GDPR}
As previously mentioned, when a \gls{sac} raises a relevant and reasoned objection to a draft decision of the \gls{lsa}, and the \gls{lsa} does not follow the objection, the matter must be referred to the dispute resolution mechanism of the \gls{edpb}.\autocite[65(1)(a)]{GDPR}
The \gls{edpb} must then issue -- within a month or two, depending on the complexity of the case -- a binding decision on the \gls{lsa} and the \glspl{sac} on the subjects of the objection.\autocite[65(2)]{GDPR} 
The \gls{lsa} is compelled to adopt a decision based on the \gls{edpb}'s decision within a month.\autocite[65(6)]{GDPR}

The \gls{edpb} can also resolve disputes between \glspl{sa} when there are conflicting views as to which \gls{sac} is competent for the main establishment of a controller or processor.\autocite[65(1)(b)]{GDPR}
Furthermore, the \gls{edpb} can function as a dispute resolution forum and issue binding decisions when an \gls{sa} refuses to comply with an opinion issued under Article 64.\autocite[65(1)(c)]{GDPR}

Finally, while as a rule, \glspl{sac} must follow the cooperation procedure\autocite[63]{facebookBelgium}, an urgent need to protect the fundamental rights and freedoms of data subject may allow an exception to that rule.\autocite[63]{facebookBelgium}
As part of the consistency mechanism, \glspl{sa} may trigger an urgency procedure, derogating from the cooperation procedure.\autocite[66(1)]{GDPR}
The urgency procedure allows any \gls{sa} to adopt urgent provisional measures, with a maximum effect of three months, on the territory of their \gls{ms}.\autocite[66(1)]{GDPR}
\glspl{sac} must inform all the other \glspl{sac}, the \gls{edpb}, and the Commission of their urgent provisional measure, and may request an urgent opinion or binding decision from the \gls{edpb}, if it considers that final measures must be urgently adopted.\autocite[66(2)]{GDPR}

The urgency procedure can also be triggered following the failure of another \gls{sa} to comply with the cooperation procedure.
When an \gls{sa} fails to comply with a request for mutual assistance, the threshold for the urgency procedure is considered to be met, allowing an \gls{sa} to adopt provisional measures, as well as requesting an urgent binding decision from the \gls{edpb}.\autocite[61(8)]{GDPR}
Additionally, when an \gls{lsa} does not invite a \gls{sac} to participate in joint operations, or ignores a request of a \gls{sac} to participate within one month, the \gls{sac} may also adopt an urgent provisional measure, requiring an urgent binding decision of the \gls{edpb}.\autocite[62(7)]{GDPR}

Finally, the \gls{edpb} can also issue urgent opinions or binding decisions on the request of an \gls{sa}, when the \gls{sa} believes there is an urgent need to act to protect the fundamental rights and freedoms of data subjects.\autocite[66(3)]{GDPR}

\subsection{The Post-Reform Cooperation and Consistency Mechanism}

As mentioned in Section 1.1.2, the Procedural Regulation only concerns cross-border processing.\autocite[1]{ProceduralRegulation} This section highlights the amended flow of cooperation between the \gls{lsa} and the \glspl{sac} and other changes in cross-border enforcement. 
For the purposes of this section, a `receiving \gls{sa}' refers to the \gls{sa} receiving the initial complaint, before the allocation of competence is determined. Additionally, the moment of confirmation of the competence of the \gls{lsa} is referred to as $T_0$.

\subsubsection{The harmonization of `parties under investigation' and complaint forms}

The Procedural Regulation provides a new definition for `parties under investigation', which refers to the controller or processor under investigation for an infringement of the \gls{gdpr} concerning cross border processing.\autocite[2]{ProceduralRegulation}

The Procedural Regulation also harmonizes the information to be contained in a complaint form, prohibiting an \gls{sa} from imposing additional requirements on the complainant.\autocite[4(1)]{ProceduralRegulation}
The harmonization only concerns complaints relating to cross-border processing.\autocite[4(1)]{ProceduralRegulation}
When a complainant does not include the required information, the receiving \gls{sa} must declare the complaint inadmissible and inform the complainant within two weeks of receipt.\autocite[4(2)]{ProceduralRegulation}

\subsubsection{The proceduralization of competence attribution}

When an \gls{sa} receives a formally valid complaint, as outlined in Section 3.3.1, it must formulate a preliminary conclusion on: (i) whether the complaint concerns cross-border processing; (ii) the \gls{sa} it considers competent as \gls{lsa}; and (iii) whether the alleged infringement exclusively affects the \gls{ms} of its own competence, as explored in section 3.2.\autocite[4(4)]{ProceduralRegulation}

The receiving \gls{sa}  must, within six weeks from receipt, transmit the complaint to the \gls{lsa}.\autocite[4(5)]{ProceduralRegulation}
The determination of the receiving \gls{sa} is binding on the \gls{lsa},\autocite[4(5)]{ProceduralRegulation} who will have six weeks to confirm its competence or contest it by referring the matter to the \gls{edpb} for dispute resolution.\autocite[4(6)]{ProceduralRegulation}
If the \gls{lsa} is inactive, at the expiration of the six weeks, the receiving \gls{sa} refers the matter to the \gls{edpb} for dispute resolution.\autocite[4(6)]{ProceduralRegulation}
The moment the competence is confirmed either by the LSA itself, or through a binding decision of the \gls{edpb} will be referred to as $T_0$.

\subsubsection{Three new tracks for cross-border enforcement cases}

The Procedural Regulation provides for three tracks for the resolution of complaints relating to cross-border processing: early resolution, simple cooperation, and complex cooperation.

\paragraph{Early resolution} Through early resolution, a complaint is dismissed when the alledged infringement has been brought to an end, making the complaint `devoid of purpose'.\autocite[5(2)]{ProceduralRegulation}
The procedure can be applied to cases that concern cross-border infringements relating to the rights of data subjects laid out in ch III of the \gls{gdpr}.\autocites[5(1)]{ProceduralRegulation}[ch III]{GDPR}
Early resolution can be triggered by the receiving \gls{sa} before it even transmits it to the \gls{lsa}, or by the \gls{lsa} after receiving the complaint.\autocite[5(1)]{ProceduralRegulation}
\todo[inline]{Reword this paragraph to include at "any time before the submission of the preliminary findings under Art. 19, or, where simple cooperation applies, before the submission of the draft decision."}


The \gls{sa} dismissing the complaint must establiosh, based on supporting evidence, that the infringement has been brought to an end.\autocite[5(2)]{ProceduralRegulation}
When a complaint is dismissed through early resolution, the dismissing \gls{sa} must inform the complainant of the end of the infringement, that it considers the complaint `devoid of purpose', the consequences of early resolution, and the complainant's right to object to early resolution within four weeks.\autocite[5(2)]{ProceduralRegulation}
When a complaint is dismissed by the receiving \gls{sa}, the receiving \gls{sa} establishes the resolution of the complaint within two weeks from the end of the objection period, and informs all parties.\autocite[5(3)]{ProceduralRegulation}
When a complaint is dismissed by the \gls{lsa}, the \gls{lsa} submits a draft decision to the \glspl{sac} within four weeks from the expiration of the objection period.\autocite[5(4)]{ProceduralRegulation}

It is important to note that when early resolution is employed, the procedure does not follow the new requirements provided for by the Procedural Regulation, including new procedural rights for complainants and parties under investigation.\autocite[5(6)]{ProceduralRegulation}

\paragraph{Simple cooperation} A complaint can be handled through a simple cooperation procedure by the \gls{lsa} when, after having formed a preliminary view on the case, the \gls{lsa} believes that: 
\begin{enumerate*}[label=(\roman*)]
    \item there is no doubt regarding the scope of the investigation;
    \item and the case does not present legal or factual complexity that would require additional cooperation with the \glspl{sac}.
\end{enumerate*}\autocite[6(1)]{ProceduralRegulation}
Simple cooperation should not be employed in cases where there is a systemic risk or recurring problems in multiple \glspl{ms}, when the case intersects with other legal fields other than data protection, or when the case affects a large number of data subjects or is based on multiple complaints across the Union.\autocite[recital 26]{ProceduralRegulation}

To trigger the simple cooperation procedure, the \gls{lsa} must inform the \glspl{sac} of its intention to handle the case through simple cooperation within six weeks from $T_0$.\autocite[6(2)]{ProceduralRegulation}
Any \gls{sac} may object to the application of the simple cooperation procedure within two weeks of the notification, and the procedure will continue according to the new complex procedure.\autocite[6(3)]{ProceduralRegulation}

Under the simple procedure, an \gls{lsa} must deliver a draft decision to the \glspl{sac} within twelve months of $T_0$.\autocite[12(6)]{ProceduralRegulation}
An LSA may extend the deadline by no more tham two months when national law requires steps to be taken by the LSA that cannot be completed within the twelve months\autocite[12(6)]{ProceduralRegulation}
Aside from that, most the new procedural requirements and time-limits established by the Procedural Regulation do not apply to the simple procedure.\autocite[6(1)]{ProceduralRegulation}
Importantly, this excludes new rights of the parties under investigation and complainants to be heard, with the simple cooperation procedure stating that the parties should have a right to be heard, but leaving it to the \glspl{ms} to determine the procedural rights of the parties.\autocite[6(4)]{ProceduralRegulation}

\paragraph{Complex cooperation}
